\documentclass[11pt,letterpaper]{article}
\usepackage[margin=1in]{geometry}
\usepackage{microtype}
\usepackage[hidelinks]{hyperref}
\usepackage{parskip}
\setlength{\parskip}{5pt}
\setlength{\parindent}{0pt}
\pagestyle{plain}
\begin{document}
\begin{center}
{\Large\bfseries AI Use Reflection}\\[5pt]
DSCI 498 Filter Surrogate Explanation Project\\
Author perspective: Guohao Xiang\quad Fall 2026
\end{center}

\textbf{Where AI participated.}
I used AI assistance to operate the ADS Python Console and automate batch and parallel circuit simulations. Repeated simulator evaluation was necessary for generating responses and comparing attribution methods against a simulator-defined reference. Automation made it practical to execute and organize many related jobs. The retained runner patches the baseline netlist, launches simulations, reads the native output, checks the frequency grid and finite values, and stores completed cases. Parallel execution also required handling failed jobs and license limits. These are workflow benefits; I did not conduct a controlled measurement of the time saved by AI.

I also used AI to help find and explain technical descriptions of the algorithms and mathematical foundations of the simulators being compared. This assistance helped me distinguish ADS component-level circuit calculations from openEMS full-wave EC-FDTD, and understand the role of Momentum as an additional full-wave comparison. The explanations supported my understanding of modeling assumptions and numerical methods. A plausible explanation of an algorithm, however, could not establish which particular mechanism caused a difference in our filter results.

\textbf{How outputs were checked.}
The project retained checks that connect the automated workflow to inspectable evidence. Known ADS responses and the complete geometry coalitions were compared with their dataset counterparts. The recorded maximum complex-response difference for the complete coalitions was zero. Such agreement supports the consistency of the circuit-simulation and data-processing pipeline for the checked cases. It does not establish that the circuit model captures every physical effect, or that every interpretation produced with AI assistance is correct.

The more consequential check was a separate full-wave evaluation in openEMS. Three geometry changes that passed the specification mask after SHAP-guided tuning in ADS were re-evaluated. None passed the same mask in the archived full-wave responses. Two cases improved substantially and one became worse. This finding changed what the project could responsibly claim: a faithful explanation and a circuit-confirmed tuning result could still inherit limitations of the model that generated them. openEMS served as an independent stress test, rather than the validation split used for training or hyperparameter selection.

\newpage
\textbf{What required judgment.}
AI-assisted source interpretation was useful for forming an account of the discrepancy, but I needed to keep the evidence boundary clear. Different solution methods, meshes, ports and boundary conditions can all matter. Agreement between full-wave calculations provides a cross-check; it does not by itself prove convergence or isolate radiation as the sole cause. The final comparison therefore distinguishes observed response changes from explanations that remain hypotheses.

Submission preparation also used AI assistance to organize documentation and recompute values from saved outputs. These checks exposed a reporting issue: model bandwidth metrics used a peak-relative 3 dB definition, while the historical full-wave summaries used an absolute minus 3 dB definition. The final report labels both rather than treating their differences as interchangeable. The old smooth-versus-hard violation error is similarly distinguished from a same-target prediction error. Preserving those qualifications matters more than presenting a single clean but misleading number.

\textbf{Responsibility and limits.}
The follow-up interview reinforced that explanations belong in the evidence layer. The practitioner wanted the frequency region, failed specification, violation magnitude and case-level quality indicators before considering a recommendation. That feedback informs presentation and permissible actions; it does not demonstrate measured gains in human decisions. I retain responsibility for the claims I submit and for distinguishing source descriptions, computed results and interpretation. In particular, local fit quality and additivity checks should not be presented as general guarantees of physical correctness.

My main lesson is to combine AI-supported execution and understanding with checks that can contradict the initial result. The automated ADS workflow enabled a broad experiment, while independent full-wave evaluation revealed a boundary that strong within-model agreement could not reveal. Future use should retain that possibility of rejection, document the exact tools and prompts more systematically, and verify engineering recommendations with the physical model or measurements appropriate to the decision.

\end{document}
